\documentclass[10pt,a4paper]{amsart}
\usepackage[margin=19mm]{geometry}
\usepackage{amsmath,amssymb,mathtools}
\usepackage[scr=rsfs,cal=euler]{mathalfa}
\usepackage{xfrac}
\usepackage[unicode,hidelinks,bookmarks=false]{hyperref}
\newcommand{\ie}{i.e.\ }
\newcommand{\cf}{cf.\ }
\newcommand{\resp}{respectively\ }
\newcommand{\Subsection}{\subsection}
\providecommand{\nobreakdash}{\nobreak\hbox{-}}
\setlength{\emergencystretch}{2em}
\multlinegap0pt
\usepackage{bm}
\usepackage{bbm}
\usepackage{dsfont}
\usepackage{xspace}
\usepackage{cancel}
\usepackage{mathtools}
\usepackage{amsthm}
\makeatletter
\@ifundefined{lemma}{\newtheorem{lemma}{Lemma}}{}
\makeatother
\title[Gromov-Hausdorff stability of global attractors of damped wa]{Gromov-Hausdorff stability of global attractors of damped wave equations under perturbations of the domain}
\author{Ngoctu Bui, Jihoon Lee}
\date{Source: arXiv:2601.13650v1 (CC BY 4.0)}
\begin{document}
\maketitle

\noindent\emph{Source and licence.} Gromov-Hausdorff stability of global attractors of damped wave equations under perturbations of the domain. Version arXiv:2601.13650v1 (CC BY 4.0), distributed under the Creative Commons Attribution 4.0 International licence, as recorded in the arXiv licence metadata for this exact version and shown on the abstract page https://arxiv.org/abs/2601.13650v1. Full rights notice: licenses/arxiv-2601.13650-CC-BY-4.0.txt.\par\smallskip

\noindent\textit{Editorial scope.} Counted unit: the Lemma whose source label is \texttt{lem in} in the article (source0.tex lines 239--249, source label \texttt{lem in}), delivered together with the article's own complete original proof of it (source0.tex lines 251--381, one \texttt{proof} environment of 131 lines). No mathematics is written, repaired or re-derived: statement and proof are the article's own text line for line. Required context retained uncounted: perturbed (lines 130--136); theo21 (lines 226--231). Every definition, display and label of those lines is retained unchanged. Omitted from this package and not redistributed: 1--129, 137--225, 232--238, 250--250, 382--597. Nothing inside a retained excerpt is omitted or altered: every retained body is delivered exactly as the article's lines are written, so no omission receipt of any kind accompanies this package. Citations: the retained excerpts carry no citation command at all, so no bibliography entry of the article is transcribed and no reference is redistributed. Reference adaptation: none was needed and none was made; every reference inside the delivered unit resolves inside it, so no unresolved reference or citation is produced. Printed slips kept literal: no printed word, symbol, hypothesis or number of the counted statement or of its proof was altered. Layout and class: the article's own class is \texttt{amsart} with packages including amssymb, latexsym, hyperref, amsmath; the wrapper is the lane's pinned \texttt{amsart} editorial wrapper at 10pt on A4, which restates the theorem-like environments the retained text needs and adds the article's own macro definitions that the retained text needs (none), with extra wrapper packages (bm, bbm, dsfont, xspace, cancel, mathtools). The delivered numbering is the unit's own. Assets: no retained span contains a figure environment, a graphics-inclusion command, a PDF-page inclusion, a table or a TikZ picture; no image file is copied into this package, the package declares no asset dependency and no image file is redistributed with it. Licence: version arXiv:2601.13650v1 of this article is distributed under the Creative Commons Attribution 4.0 International licence, as recorded in the arXiv licence metadata for this exact version and shown on the abstract page https://arxiv.org/abs/2601.13650v1. A full rights notice is delivered at licenses/arxiv-2601.13650-CC-BY-4.0.txt. Characters: every retained excerpt is pure ASCII, so the delivered engine, XeLaTeX with its default Latin Modern text font, reports no missing character.
\begin{equation} \label{perturbed}
\begin{cases}
u_{tt} + u_t - \Delta u = -f(u)  ~& \mbox{in}\quad \Omega_\delta \times (0,\infty), \\
u=0 &\mbox{on} \quad \partial \Omega_\delta \times (0, \infty), \\
u(0,x) = u_{0,\delta}(x), \quad u_t(0,x)=u_{1,\delta}(x) & \mbox{in}  \quad \Omega_\delta, 
\end{cases}
\end{equation}

\begin{lemma}\label{theo21}
    Let $0\leq \delta\leq 1$. There exist a positive constant $C>0$ and, for any $r_0>0$, $r_1>0$, two positive constants $C^*(r_0)$ and $C^*(r_1)$ such that, for any solution $u_\delta(t)$ of equation \eqref{perturbed} with $\|u_{1,\delta}\|^2_i + \|u_{0,\delta}\|^2_{i+1} \leq r_i ^2, (i=0,1),$ we have the following estimate: for any $t\geq 0$,
    \begin{align*}
        \left\|\frac{\partial^2 u_\delta (t)}{\partial^2 t}\right\|^2_0 +\left\|\frac{\partial u_\delta (t)}{\partial t}\right\|^2_1 + \|u_\delta(t)\|^2_2 \leq C^*(r_0)+ C^*(r_1)e^{-C t}.
    \end{align*}
\end{lemma}

\begin{lemma}\label{lem in}
    Let $\{h_n\}_{n\in \mathbb{N}}$ be a sequence in $\mathcal P(\Omega)$, and $\{V_{0,n}\}_{n\in \mathbb{N}}$ be a sequence in $X^0_{\delta_0}$  which is bounded in $X^1_{\delta_0}$ such that
    \begin{align*}
        h_n \rightarrow h_0 \mbox{ and } V_{0,n}\rightarrow V_0 \mbox{ as } n\rightarrow \infty.
    \end{align*}
    Then for any $T>0$, we have
    \begin{align*}
        \|I_n(T_{\delta_n}(I_n^{-1}(V_{0,n}),t))-T_{\delta_0}(V_0, t)\|_{X^0_{\delta_0}}\rightarrow 0 \mbox{ for any } t\in [0, T],
    \end{align*}
    where $\delta_n = d_{C^2}(h_n, 1_\Omega)$ for $n = 0,1,2,...$. 
\end{lemma}

\begin{proof}
    Take $T>0$, and let $U_{0,n} = I_n^{-1}(V_{0,n})=(u_{0,n},u_{1,n})$. Let  $u_n(t)$ be the weak solution of the equation
    \begin{align}\label{eq:3.1}
        \left\{\begin{array}{ll}
u_{tt} + u_t - \Delta u = -f(u) ~& \mbox{in}\quad \Omega_{\delta_n} \times (0,\infty), \\
u=0 &\mbox{on} \quad \partial \Omega_{\delta_n} \times (0, \infty), \\
u(0,x) = u_{0,n}(x), \quad u_t(0,x)=u_{1,n}(x),& \mbox{in}  \quad \Omega_{\delta_n},
\end{array}\right.\end{align}
and $v(t)$  be the weak solution of the equation
\begin{align}\label{eq:3.2}
    \left\{\begin{array}{ll}
v_{tt} + v_t - \Delta v = -f(v)~& \mbox{in}\quad \Omega_{\delta_0} \times (0,\infty), \\
v=0 &\mbox{on} \quad \partial \Omega_{\delta_0} \times (0, \infty), \\
v(0,x) = v_0(x), \quad v_t(0,x)=v_1(x),& \mbox{in}  \quad \Omega_{\delta_0},
\end{array}\right.\end{align}
Let $U_n(t) = (u_n(t), \partial_t u_n(t))$, $V(t) = (v(t), \partial_t v(t))$, $V_n = I_n(U_n)$ and $Z_n = V_n - V$. We need to show that $\|Z_n (t)\|_{X^0_{\delta_0}}\rightarrow 0 \mbox{ for all } t\in [0, T].$ 

Since $V_{0,n} $ is bounded in $X^1_{\delta_0}$ for $n=0,1,2,...$, and $I_n^{-1} $ is a bounded isomorphism from $X_{\delta_0}^i$ to $X_{\delta_n}^i$ for $i\in \{0,1\}$, we have that the initial values $U_{0,n}$ is also bounded in $X_{\delta_n}^1$. Hence, by Lemma \ref{theo21}, we obtain the boundedness of $U_n(t)$ in $X_{\delta_n}^1$. Similarly we see that $V_n(t)$ belongs to the space $X_{\delta_0}^1$. Hence we know that $Z_n \in X_{\delta_0}^1$, and so $i_n^{-1}(\partial_t z_n) \in L^2(0,T;H_0^1(\Omega_{\delta_n}))$. We take the inner product with $i_n^{-1}(\partial_t z_n)$ on the both sides of equation \eqref{eq:3.1} to deduce that
\begin{align}\label{eq:3.3}
\int_{\Omega_{\delta_n}}\partial_{tt}u_n (i_n^{-1}(\partial_t z_n)) dx &+ \int_{\Omega_{\delta_n}}\partial_tu_n (i_n^{-1}(\partial_t z_n)) dx \notag\\
&+\int_{\Omega_{\delta_n}}\nabla u_n \cdot \nabla (i_n^{-1}(\partial_tz_n))dx 
=-\int_{\Omega_{\delta_n}}f(u_n) i_n^{-1}(\partial_t z_n) dx.
\end{align}
Let $y = h_0 \circ h_n^{-1}(x)$ for $x \in \Omega_{\delta_n}$. By the change of variable and \eqref{eq:3.3} (using the notation $x$ instead of $y$ for convenience), we get  that
\begin{align} \label{eq:3.4}
&\int_{\Omega_{\delta_0}}\partial_{tt}v_n + \partial_t z_n|\det H_n|dx+ \int_{\Omega_{\delta_0}}\partial_{t}v_n  \partial_t z_n|\det H_n|dx\notag \\
&+\int_{\Omega_{\delta_0}}\overline{H}_n\nabla v_n \cdot \overline{H}_n\nabla \partial_t z_n |\det H_n| dx =-\int_{\Omega_{\delta_0}} f(v_n) \partial_t z_n|\det H_n| dx,
\end{align}
where
\begin{equation*}\label{eq:3.5}
H_n(x)=D(h_n\circ h_0^{-1}(x))=\left(\dfrac{\partial (h_n\circ h_0^{-1}(x))_i}{\partial x_j}\right)_{N\times N}, \; x\in \Omega_{\delta_0},
\end{equation*}
and $\overline{H}_n=(H_n^{-1})^{\mathsf{T}}$ is the transpose of the inverse $H_n^{-1}$.
\vskip 5 mm
We take the inner product with $\partial_t z_n$ on both sides of the equation \eqref{eq:3.2} to deduce that
\begin{equation}\label{eq:3.6}
\int_{\Omega_{\delta_0}}v_{tt}\partial_t z_n dx+\int_{\Omega_{\delta_0}}v_t \partial_t z_n dx+ \int_{\Omega_{\delta_0}}\nabla v\cdot \nabla \partial_t z_ndx=-\int_{\Omega_{\delta_0}}f(v) \partial_t z_ndx.\end{equation}
By subtracting \eqref{eq:3.6} from \eqref{eq:3.4}, we get
\begin{align}\label{eq:3.7}
    &\int_{\Omega_{\delta_0}}\partial_{tt}z_n \partial_tz_n |\det H_n| dx+\int_{\Omega_{\delta_0}}\partial_{t}z_n \partial_tz_n |\det H_n| dx\\
&\;\;\;+\int_{\Omega_{\delta_0}}(|\det H_n|-1)(v_t+v_{tt}) \partial_tz_n dx +\int_{\Omega_{\delta_0}}\overline{H}_n\nabla z_n \cdot \overline{H}_n\nabla \partial_t z_n |\det H_n| dx\notag\\
&=\int_{\Omega_{\delta_0}}\left(I-\overline{H}_n\right)\nabla v \cdot \overline{H}_n\nabla \partial_tz_n |\det H_n| dx\notag\\
&\;\;\;+\int_{\Omega_{\delta_0}}\nabla v \cdot (I-\overline{H}_n)\nabla \partial_tz_n |\det H_n|dx+\int_{\Omega_{\delta_0}}(1-|\det H_n|)\nabla v\cdot \nabla \partial_tz_ndx\notag\\
&\;\;\;+\int_{\Omega_{\delta_0}}(1-|\det H_n|)f(v) \partial_tz_ndx-\int_{\Omega_{\delta_0}} (f(v_n)-f(v)) \partial_tz_n|\det H_n| dx\notag.
\end{align}

Now we estimate each term. For the left hand side of \eqref{eq:3.7}, the first term is
\begin{align*}
    \int_{\Omega_{\delta_0}}\partial_{tt}z_n \partial_tz_n |\det H_n| dx = \frac{1}{2} \frac{d}{dt}\int_{\Omega_{\delta_0}}\left|\partial_t z_n\right|^2 |\det H_n| dx,
\end{align*}
the second term is
\begin{align*}
    \int_{\Omega_{\delta_0}}\partial_t z_n \partial_tz_n |\det H_n| dx = \int_{\Omega_{\delta_0}}\left|\partial_t z_n\right|^2 |\det H_n| dx,
\end{align*}
the third term is
\begin{align*}
    - \int_{\Omega_{\delta_0}}(|\det H_n|-1)(v_t+v_{tt}) \partial_tz_n dx
    &\leq \||\det H_n| -1\|_{\infty} \int_{\Omega_{\delta_0}}\left(- \nabla v\cdot \nabla z_n - f(v) z_n \right)dx\\
    &\leq \frac{1}{2}\||\det H_n| -1\|_{\infty} K(t),
\end{align*}
where $K(t) = \|v\|_1^2 + \|\partial_t z_n\|^2_1 +\|f(v)\|^2_0 +\|\partial_t z_n\|^2_0$, 
and the fourth term is
\begin{align*}
    \int_{\Omega_{\delta_0}}\overline{H}_n\nabla z_n \cdot \overline{H}_n\nabla \partial_t z_n |\det H_n| dx = \frac{1}{2} \frac{d}{dt}\int_{\Omega_{\delta_0}}|\overline{H}_n\nabla z_n|^2|\det H_n|dx.
\end{align*}

For the right hand side of \eqref{eq:3.7}, we have
\begin{align*}
    &\int_{\Omega_{\delta_0}}\left(I-\overline{H}_n\right)\nabla v \cdot \overline{H}_n\nabla \partial_tz_n |\det H_n| dx \\&\;\;\;+\int_{\Omega_{\delta_0}}\nabla v \cdot (I-\overline{H}_n)\nabla  \partial_tz_n |\det H_n|dx+\int_{\Omega_{\delta_0}}(1-|\det H_n|)\nabla v\cdot \nabla \partial_tz_ndx \\
    &\leq \dfrac{1}{2}\Big\{\|I-\overline{H}_n\|_{\infty}\|\det H_n\|_{\infty}\left(\|\overline{H}_n\|_{\infty}+1\right) +\||\det H_n |-1\|_{\infty}\Big\}\left(\|v\|_1^2+\|\partial_t z_n\|_1^2\right),
\end{align*}
the next term can be estimated as
\begin{align*}
    \int_{\Omega_{\delta_0}}(1-|\det H_n|)f(v) \partial_t z_ndx & \leq \||\det H_n |-1\|_{\infty}\|f(v)\|_0\|\partial_tz_n\|_0\\
&\leq \||\det H_n |-1\|_{\infty}\left(\dfrac{1}{2}\|f(v)\|^2_0+\dfrac{1}{2}\|\partial_tz_n\|^2_0\right),
\end{align*}
and the last term can be estimated by
\begin{align*}
    &-\int_{\Omega_{\delta_0}} (f(v_n)-f(v)) \partial_t z_n|\det H_n| dx \leq l \int_{\Omega_{\delta_0}}|z_n \partial_t z_n||\det H_n|dx \\
    &\leq \frac{l}{2}\|z_n \sqrt{|\det H_n|}\|^2_0 +\frac{l}{2}\|\partial_t z_n \sqrt{|\det H_n|}\|^2_0 \leq \frac{\ell}{2}(\|\nabla_{\delta_0} z_n \sqrt{|\det H_n|}\|^2_0 +\|\partial_t z_n \sqrt{|\det H_n|}\|^2_0)\\
    &\leq \ell(\|\overline{H}_n \nabla_{\delta_0}z_n \sqrt{|\det H_n|}\|^2_0+\|\partial_t z_n \sqrt{|\det H_n|}\|^2_0)+\ell \|I - \overline{H}_n\|^2_\infty\|\det H_n\|_{\infty}\|\nabla_{\delta_0}z_n \|_0^2,
\end{align*}
where $\ell = \max\left\{\frac{l}{\lambda_1}, l\right\}$, and $\lambda_1$ is the first eigenvalue of the operator $A_0$. 


Consequently we derive
\begin{align*}
    &\frac{1}{2} \frac{d}{dt}\int_{\Omega_{\delta_0}}\left|\partial_t z_n\right|^2 |\det H_n| dx + \frac{1}{2} \frac{d}{dt}\int_{\Omega_{\delta_0}}|\overline{H}_n\nabla z_n|^2|\det H_n|dx \\
    &\leq \||\det H_n| -1\|_{\infty} K(t) \\
    &\;\;\;+\dfrac{1}{2}\Big\{\|I-\overline{H}_n\|_{\infty}\|\det H_n\|_{\infty}\left(\|\overline{H}_n\|_{\infty}+1\right) +\||\det H_n |-1\|_{\infty}\Big\}\left(\|v\|_1^2+\|\partial_t z_n\|_1^2\right)\\
    &\;\;\;+ \ell(\|\overline{H}_n \nabla z_n \sqrt{|\det H_n|}\|^2_0+\|\partial_t z_n \sqrt{|\det H_n|}\|^2_0)\\ 
    &\;\;\;+\ell\|I - \overline{H}_n\|^2_\infty\|\det H_n\|_{\infty}\|\nabla z_n \|_0^2, 
\end{align*}
which is equivalent to
\begin{align}\label{eq:3.11}
    &\frac{d}{dt}\left(e^{-2\ell t}\int_{\Omega_{\delta_0}}\left(\left|\partial_t z_n\right|^2 +|\overline{H}_n\nabla_{\delta_0} z_n|^2\right)|\det H_n|dx\right)\\
    &\leq 2 e^{-2\ell t} \||\det H_n| -1\|_{\infty} K(t) \notag\\
    & \;\;\;+ e^{-2\ell t} \Big\{\|I-\overline{H}_n\|_{\infty}\|\det H_n\|_{\infty}\left(\|\overline{H}_n\|_{\infty}+1\right)\big\}\left(\|v\|_1^2+\|\partial_t z_n\|_1^2\right) \notag \\
    &\;\;\; + e^{-2\ell t}||\det H_n |-1\|_{\infty}\left(\|v\|_1^2+\|\partial_t z_n\|_1^2\right)\notag\\
    &\;\;\;+ e^{-2\ell t} \ell\|I - \overline{H}_n\|^2_\infty\|\det H_n\|_{\infty}\|\nabla_{\delta_0}z_n \|_0^2\notag.
\end{align}
By integrating \eqref{eq:3.11} from $0$ to $t\in [0,T]$, we derive that
\begin{align}\label{eq:3.12}
& \int_{\Omega_{\delta_0}}\left(\left|\partial_t z_n\right|^2 +|\overline{H}_n\nabla_{\delta_0} z_n|^2\right)|\det H_n|dx\\
&\leq e^{2\ell t}\int_{\Omega_{\delta_0}}\left(\left|\partial_t z_n(0)\right|^2 +|\overline{H}_n\nabla_{\delta_0} z_n(0)|^2\right)|\det H_n|dx \notag \\
&\;\;\ + 2e^{2\ell t}\||\det H_n |-1\|_{\infty}\int_0^T K(t)dt \notag \\
& \;\;\;+ e^{2\ell t} \ell\|I - \overline{H}_n\|^2_\infty \|\det H_n\|_\infty\|z_n \|^2_{L^{\infty}(0, T, H^1_0(\Omega_{\delta_0}))}\notag\\
&\;\;\;+e^{2\ell t}\Big\{\|I-\overline{H}_n\|_{\infty}\|\det H_n\|_{\infty}\left(\|\overline{H}_n\|_{\infty}+1\right)+\||\det H_n |-1\|_{\infty}\Big\} \notag\\
&\;\;\;\times  \int_0^T\left(\|v(t)\|_1^2+\|\partial_t z_n(t)\|_1^2\right)dt\notag.
\end{align}

We now choose a constant $C$ independent of $n$ so that 
\begin{align*}
\max \big\{ \int_0^TK(t)dt,~\int_0^T\left(\|v(t)\|_1^2+\|\partial_t z_n(t)\|^2_1\right)dt,~ \|z_n \|^2_{L^{\infty}(0, T, H^1_0(\Omega_{\delta_0}))}\big\} \le C.
\end{align*}
Since $h_n \to h_0$ in $\mathcal{P}(\Omega)$ and $\det$ is a continuous operator, we can see that  $\|\det H_n\|_{\infty}$  and  $\|\overline{H}_n\|_{\infty} $ are uniformly bounded with respect to $n$. Moreover, note that 
$$
\inf\limits_{x \in \Omega_{h_0}}|\det H_n(x)| \ge \frac12 ~\mbox{and}~\inf\limits_{x \in \Omega_{h_0}} |\overline{H}_n(x)| \ge \frac12
$$ for sufficiently large $n \in \mathbb N$, and 
\begin{equation*}
\||\det H_n |-1\|_{\infty}\to 0\mbox{ and }   \|I-\overline{H}_n\|_{\infty}\to 0 \mbox{ as } n\to \infty.
\end{equation*}
Therefore,  from \eqref{eq:3.12} we conclude that
\begin{align*}
    &\frac{1}{4}\int_{\Omega_{\delta_0}}\left(\left|\partial_t z_n\right|^2 +|\nabla z_n|^2\right)dx\\
    &\leq \inf_{x\in \Omega_{\delta_0}}|\det H_n(x)|
    \int_{\Omega_{\delta_0}}\left(\left|\partial_t z_n\right|^2 +|\overline{H}_n(x)\nabla z_n|^2\right)dx\\
    &\leq e^{2\ell T} \|\det H_n\|_{\infty}\int_{\Omega_{\delta_0}}\left(\left|\partial_t z_n(0)\right|^2 +\|\overline{H}_n\|^2_\infty|\nabla z_n(0)|^2\right)dx \to 0 \mbox{ as } n\to \infty,
\end{align*}
which completes the proof of the lemma.
\end{proof}

\end{document}
